\chapter{绪论}

% ===== 提示 =====
% 绪论是"提出问题"的部分，一般主要概述：
% 1. 课题研究问题背景及研究价值
% 2. 国内外研究现状
% 3. 主要研究内容、方法或设计思路
% 总字数控制在2000字左右。
% ================

% ===================================================================
% ★ 参考文献引用完整指南
% ===================================================================
%
% 【第一步】在 references.bib 中添加文献条目
%   每条文献有一个唯一的 key（如 example_book），格式参考 bib 文件中的示例。
%
% 【第二步】在正文中引用
%   \cite{key}            → 生成编号，如 [1]
%   \cite{key1, key2}     → 同时引用多篇，如 [1,2]
%
% 【引用位置说明】
%   - 引用放在句末标点之前，如：……具有重要意义\cite{key}。
%   - 只有在 references.bib 中存在 且 在正文中被 \cite{} 引用过的文献，
%     才会出现在最终的"参考文献"列表中。
%   - 参考文献按引用顺序自动编号，无需手动排序。
%
% 【编译说明】
%   引用和参考文献需要多次编译才能正确显示：
%   xelatex → biber → xelatex → xelatex（compile.bat 脚本已自动完成）
%   直接运行 compile.bat（Windows）或 bash compile.sh（Linux/macOS）即可，
%   无需手动多次编译。
%
% 【常见问题】
%   - 引用显示为 [?]：说明 biber 未运行或 key 拼写错误，请运行完整编译脚本
%   - 参考文献不少于10条（学校规范要求）
%   - 应包含专著、论文、报刊等多种类型，其中应有外文参考文献
% ===================================================================

\section{研究背景与意义}

在此撰写你的研究背景和选题意义。可以从社会需求、行业现状、技术发展趋势等方面展开论述，说明为什么选择这个课题进行研究\cite{example_book}。

\section{国内外研究现状}

\subsection{国外研究现状}

在此概述国外在该领域的研究进展和代表性成果。引用文献时在句末标点前使用\verb|\cite{}|命令，例如\cite{example_article}。同时引用多篇文献写法：\verb|\cite{key1, key2}|，例如\cite{example_article, example_conference}。

\subsection{国内研究现状}

在此概述国内在该领域的研究进展和代表性成果\cite{example_thesis}。

\section{主要研究内容}

简要说明本论文的主要研究内容和章节安排：

第1章为绪论，介绍了研究背景、国内外研究现状和主要研究内容。

第2章为相关技术介绍，对论文涉及的核心技术进行阐述。

第3章为系统设计与实现，详细描述系统的总体设计方案和具体实现过程。

第4章为总结与展望，对全文进行总结并展望未来的研究方向。
